\documentclass[10pt,a4paper]{article}
\usepackage[margin=22mm]{geometry}
\usepackage{amsmath,amssymb,booktabs,array}
\usepackage[colorlinks=true,linkcolor=blue,urlcolor=blue]{hyperref}
\usepackage{microtype}
\title{RM9 continuation: priority allocation, control effects\\and conservative lane-flow readiness}
\author{Method-first ramp-metering review}
\date{1 October 2026}
\begin{document}
\maketitle
\noindent\textbf{Disposition.} RM9-UG1.A is approved as \textbf{Default}.
The selected merge subproblem is transcribed below. The complete lane-flow
coupling is not adopted; Section~\ref{sec:proposal} is an explicit proposal for
RM9-UG2. No traffic simulation, numerical derivative check or useful-gradient
certificate is reported. Proofs and examples are reviewer derivations.

\section{The approved two-input allocator}
Let $d_m,d_r\geq0$ be local mainline and ramp sending offers and $S\geq0$
their common receiving rate, all in vehicle equivalents per unit time.
Let structural priorities $p_m,p_r>0$ be independent of offers and meter
rate, with $a=p_m/(p_m+p_r)\in(0,1)$. A transcription of [1]'s MISO law is
\begin{align}
 r&=\min\{d_r,\max\{(1-a)S,S-d_m\}\},\label{eq:r}\\
 m&=\min\{d_m,S-r\}.\label{eq:m}
\end{align}
The max term lets the ramp use supply the mainline cannot consume; the
second line lets the mainline consume supply the ramp cannot consume.
It is equivalent to the symmetric expression
$m=\min\{d_m,\max\{aS,S-d_r\}\}$.
This is accepted physical flow, not an optimized assignment or guaranteed entry.

\paragraph{Bounds and work conservation.}
The inner max lies in $[0,S]$, hence $0\leq r\leq\min(d_r,S)$ and
$0\leq m\leq d_m$, $m+r\leq S$. If $d_m+d_r\leq S$, then
$S-d_m\geq d_r$, giving $r=d_r,m=d_m$.
If $d_m+d_r>S$ and $d_m\geq aS$, then $r=\min(d_r,(1-a)S)$:
when the ramp is under its share, $m=S-d_r$; otherwise $m=aS$.
If $d_m<aS$, then $r=S-d_m,m=d_m$. Consequently
\begin{equation}
 m+r=\min\{S,d_m+d_r\}.\label{eq:total}
\end{equation}
This proof covers zero supply/offers. No source-wide MIMO/FIFO behavior follows.

\paragraph{Ordinary derivatives, not a verified whole-chain gradient.}
At fixed $d_m,S,a$, write $B=\max((1-a)S,S-d_m)$.
Away from $d_r=B$,
\begin{equation}
 \frac{\partial r}{\partial d_r}=\begin{cases}1,&d_r<B,\\0,&d_r>B,\end{cases}
 \qquad
 \frac{\partial(m+r)}{\partial d_r}
 =\begin{cases}1,&d_m+d_r<S,\\0,&d_m+d_r>S.\end{cases}
\end{equation}
At ties the ordinary derivative need not exist; no smoothing or generalized
estimator is adopted. In a supply-limited region, increasing ramp service can
reduce mainline service one for one. This does not prove a useful meter gradient.
\newpage

\section{What the rule changes in the control chain}
\subsection*{Hand-calculated service examples}
Use $S=180$ veh-equivalents/h and $a=2/3$ solely for illustration.
The nominal mainline/ramp shares are 120/60. No site calibration is implied.
\begin{center}
\begin{tabular}{rrrrl}
\toprule
$d_m$ & $d_r$ & $m$ & $r$ & Total / interpretation\\
\midrule
180&30&150&30&180; mainline uses spare ramp share\\
180&10&170&10&180; allocation changes, total unchanged\\
180&80&120&60&180; 20 of ramp offer remains unserved\\
80&70&80&70&150; all eligible offers served\\
80&40&80&40&120; less offer leaves more idle supply\\
\bottomrule
\end{tabular}
\end{center}
Rows compare local offers, not instantaneous upstream meter values. An upstream
rate $u(t)$ bounds release $e(t)$; stored/traveling ramp traffic determines
$d_r(t)$ later. A co-located meter is a separate geometric specialization.
Post-meter traffic and blocked/unserved demand must remain in accounting.

\subsection*{Spare supply on another lane is not freely transferable}
For two separate receivers of supply 180 each, ignore lateral access at the
instant considered. Offers $(d_{m1},d_{m2},d_r)=(180,50,30)$ give
\begin{equation}
 F=\min(180,d_{m1}+d_r)+\min(180,d_{m2})=180+50=230.
\end{equation}
Offers $(90,140,30)$ have the same total offer 260 but give $F=120+140=260$.
Thus lane-specific eligibility can matter to total service. These are offer-level
examples, \textbf{not matched density/flow/speed histories}, feasible controlled
trajectories, or a proof that lane contrast changes the best meter action.
A physical redistribution law is needed to connect them to traffic dynamics.

\subsection*{Why queue redistribution may produce no objective benefit}
With all affected uncompleted mass $M$, common arrivals $A$, final system exit
rate $E$, and common clearance horizon $T_c$,
\begin{equation}
 \dot M=A-E,\qquad J=\int_0^{T_c}M(t)\,dt,\qquad M(T_c)=0.
\end{equation}
Identical initial total mass, arrivals and exit history imply identical $J$,
even if ramp and mainline queues differ. The local merge flow $m+r$ is not
necessarily $E$: downstream travel, storage and bottlenecks must be included.
On a fixed differentiable branch, for common forcing and a control parameter $b$,
\begin{equation}
 \frac{dJ}{db}=-\int_0^{T_c}(T_c-t)\frac{\partial E(t;b)}{\partial b}\,dt.
\end{equation}
Allocation changes next lane/ramp masses and may change later receiving,
spillback, lane access or exit timing. Those are possible causal paths, not
established effects in this task. If final departure history stays unchanged,
queue relocation alone gives zero total-time improvement.
\newpage

\section{Why the selected allocator does not close the lane model}
\subsection*{A source coefficient range diagnostic}
The fuller source [2], equation 6, uses a coefficient whose local numerator
contains a combined incentive. Let $x>0,y\geq0$ denote donor and target
weighted densities on a constant profile, and denote that coefficient by $P$.
From the printed expression, reviewer algebra gives
\begin{equation}
 P>1\quad\Longleftrightarrow\quad (I-1)x>2y
\end{equation}
on its positive branch. For the printed free-flow cooperation incentive near
a merge, with base incentive one and other incentives inactive, $I=2+y/x$.
Thus the isolated coefficient is $2x/(x+y)$, exceeding one if $x>y$.
A coefficient called a fraction is therefore not automatically a donor-bounded
accepted transfer. This diagnosis does \emph{not} establish reachability of every
such combination or failure of every complete source trajectory. Later operators
or domain restrictions might resolve it; they must be shown, not assumed.

\subsection*{Zero-receiving special branch}
For capacity $C>0$, supply $S\in[0,C]$ and positive proposed lateral rate $\phi$,
source equations 16--17 give $\vartheta=1-S/C$ and the printed accepted rate
\begin{equation}
 \lambda=\begin{cases}\phi,&\vartheta=1,\\
 (1-\vartheta)\phi,&\vartheta<1.\end{cases}
\end{equation}
At $S=0$ this isolated operator accepts $\phi$, whereas its limit from $S>0$
is zero. It cannot be used as a simple receiving-bounded interface without
an additional interpretation/domain condition. No branch is silently corrected.
Source equations 20--22 also modify demand and supply after lateral transfer;
applying another independent node operator could double-count budgets.
This is an unresolved joint-update transcription, not a blanket source rejection.

\subsection*{An inspected predecessor does not supply the missing autonomy}
Source [3]'s 2014 formulation optimizes lateral, longitudinal and ramp flows
jointly. Its constraints are useful conservation precedents, but its lateral
flow is not autonomous response to a meter-only action. The 2015 Part I article
[4] has verified primary metadata; its full methods were not retrieved in this
pass. A bounded attractiveness equation reproduced in [2], equation 5, is an
inspected antecedent, not verification of [4]'s complete model.

\subsection*{What has and has not passed}
\begin{center}
\begin{tabular}{>{\raggedright\arraybackslash}p{0.46\linewidth}>{\raggedright\arraybackslash}p{0.44\linewidth}}
\toprule
Passed in this analytic review & Still required\\
\midrule
Two-input allocation bounds, work conservation and fixed-branch derivatives
& Complete conservative lateral/longitudinal/ramp update\\
Local service examples and full-time accounting identity
& Exact observer/initialization and feasible identical-history pair\\
Source range/branch and actuator-scope diagnostics
& Action relevance, full derivative correctness and independent usefulness\\
\bottomrule
\end{tabular}
\end{center}
\newpage

\section{A concrete conservative adaptation for RM9-UG2.A}\label{sec:proposal}
\noindent\textbf{Proposed, not adopted.} Use two physical mainline lanes per
cell, with no instantaneous rest-group mixing. The update below is a reviewer
adaptation; it is not the complete calibrated algorithm in [2] or [4].
It retains only discretionary density-driven transfer, ordinary first-order
receiving and the Default ramp allocator. It adds no capacity-drop penalty.

For cell length $L$, timestep $\Delta t$, lane mass $n_\ell=Lk_\ell$,
maximum mass $N_\ell=Lk_\ell^{\rm jam}$ and known positive parameters,
\begin{equation}
 D_\ell(n)=\min(v_\ell n/L,C_\ell),\qquad
 S_\ell(n)=\min(C_\ell,w_\ell(N_\ell-n)/L).
\end{equation}
Use the step restrictions $\Delta t v_\ell/L\leq1$ and
$\Delta t w_\ell/L\leq1$. For $0\leq n\leq N_\ell$ they imply
$\Delta t D_\ell(n)\leq n$ and $\Delta t S_\ell(n)\leq N_\ell-n$.
Parameters may differ between lanes; symmetry is not imposed to hide state.

\paragraph{Stage 1: bounded autonomous within-cell lateral exchange.}
For structural $\mu_{\ell j}\in[0,1]$, $\kappa_{\ell j}\in(0,1]$, use the
bounded attractiveness form reproduced in [2], equation 5:
\begin{align}
 A_{\ell j}&=\mu_{\ell j}\max\left(0,
 \frac{\kappa_{\ell j}k_\ell-k_j}{\kappa_{\ell j}k_\ell+k_j}\right),\\
 \lambda_{\ell j}&=\min(A_{\ell j}D_\ell(n_\ell),S_j(n_j)),\\
 n_\ell^L&=n_\ell+\Delta t(\lambda_{j\ell}-\lambda_{\ell j}).
\end{align}
At zero denominator set $A_{\ell j}=0$ (both lane masses zero). Otherwise
$0\leq A_{\ell j}\leq1$. There is at most one adjacent target per donor in
the two-lane case. Both directions cannot be positive: that would require
$k_j<\kappa_{\ell j}k_\ell$ and $k_\ell<\kappa_{j\ell}k_j$, contrary to
$\kappa_{\ell j}\kappa_{j\ell}\leq1$. The bounds imply
\begin{equation}
 0\leq n_\ell^L\leq N_\ell,\qquad
 n_1^L+n_2^L=n_1+n_2.
\end{equation}
Rejected propensity does not leave the donor. For more lanes the same formula
cannot be applied independently without a joint donor/receiver budget.

\paragraph{Stage 2: recompute state limits and retain unspent shared budgets.}
Recomputing from $n^L$ alone could spend the original receiving or sending
budget twice. Define residual longitudinal budgets for each cell-lane:
\begin{align}
 \overline D_{i\ell}&=\min\{D_{i\ell}(n_{i\ell}^L),
 D_{i\ell}(n_{i\ell})-\lambda_{\ell j}\},\\
 \overline S_{i\ell}&=\min\{S_{i\ell}(n_{i\ell}^L),
 S_{i\ell}(n_{i\ell})-\lambda_{j\ell}\}.
\end{align}
These residuals are nonnegative by Stage 1's bounds. At an ordinary interface
use $q_{i\ell}=\min(\overline D_{i\ell},\overline S_{i+1,\ell})$.
At the shared shoulder receiver apply \eqref{eq:r}--\eqref{eq:m} with the
mainline residual donor offer and receiver residual $\overline S$.
With total accepted longitudinal/ramp incoming rate $f_{i\ell}^{\rm in}$,
\begin{equation}
 n_{i\ell}^{+}=n_{i\ell}^L+
 \Delta t(f_{i\ell}^{\rm in}-q_{i\ell}),\qquad
 f_{i\ell}^{\rm in}\leq\overline S_{i\ell},\quad
 q_{i\ell}\leq\overline D_{i\ell}.
\end{equation}
Thus incoming lateral plus longitudinal/ramp flows share the original $S$,
and outgoing lateral plus longitudinal flows share the original $D$.
Post-lateral limits also imply $0\leq n_{i\ell}^{+}\leq N_{i\ell}$.
Every internal transfer cancels between its donor/receiver. Ramp/external
waiting and post-meter propagation need corresponding conservative balances.
This is an explicitly conservative shared-budget adaptation, not calibration,
a selected geometry, a history certificate or verified useful gradient.
Stage order and step dependence can restrict flows and must remain explicit.
\newpage

\section{Next decision and evidence order}
\textbf{RM9-UG2} asks whether to adopt the restricted candidate above for
analytic review, or retain the fuller incentive/cooperation family and first
resolve its source-level coupling. The expanded alternatives, pros, cons and
dependencies are in \texttt{reporting/rm9\_30\_allocator\_contract\_and\_control\_chain.md}.
The provisional recommendation is the restricted candidate, with no positive
outcome assumed. It can legitimately establish aggregate sufficiency or no
useful metering distinction. Fuller cooperation/capacity-loss mechanisms remain
possible extensions; their necessity is not yet proved.

After adoption of a defensible law, transcribe the actual strong density,
flow, speed, cumulative-count and ramp-history functions with temporal/spatial
support, initialization, common forcing and known parameters. Check reconstruction
before seeking an exact feasible hidden-history pair. A scalar equivalent of
lane contrast remains admissible. Do not use equal speeds, weakened speed
observations, hidden prehistory or an imposed service penalty to produce ambiguity.

Only after forward action relevance should the review claim a candidate
sensitivity path. Ordinary branch derivatives are insufficient at switching
boundaries, and they do not verify independent downstream usefulness. No
finite-difference/autograd run, optimizer, SUMO experiment or successor is
opened. RM9 remains open at the physical-law decision; RM9-UG1 stays closed.

\section*{Primary references and provenance}
\begin{enumerate}
\item Wright et al. (2017), \emph{On node models for high-dimensional road
networks}. Existing RM0018, published author-hosted version, Section 3.1.
\href{https://doi.org/10.1016/j.trb.2017.09.001}{DOI 10.1016/j.trb.2017.09.001}.
\item Nagalur Subraveti, Knoop and van Arem, \emph{First order multi-lane
traffic flow model---an incentive based macroscopic model to represent lane
change dynamics}. Existing FB0030, final PDF, printed pp. 1762--1766.
\href{https://repository.tudelft.nl/file/File_2c4359e7-5e5a-4651-964f-6f3610ab89ae}{Primary final PDF}.
\item Roncoli, Papageorgiou and Papamichail (2014), \emph{Optimal control for
multi-lane motorways in presence of vehicle automation and communication systems}.
Existing FB0034, IFAC proceedings pp. 4179--4181.
\href{https://skoge.folk.ntnu.no/prost/proceedings/ifac2014/media/files/1657.pdf}{Primary proceedings PDF}.
\item Roncoli, Papageorgiou and Papamichail (2015), \emph{Traffic flow optimisation
in presence of vehicle automation and communication systems---Part I}.
RM9-L01 \textbf{metadata lead only}; no full-method adoption.
\href{https://doi.org/10.1016/j.trc.2015.06.014}{DOI 10.1016/j.trc.2015.06.014}.
\end{enumerate}
\noindent\textbf{Ownership.} The allocator transcription, analytical proofs,
hand-calculated illustrations, range/branch diagnostics and proposed split
update are reviewer work under their stated assumptions. They are not findings
reported by all the cited papers or a verified combined literature gradient chain.
\end{document}
